\begin{table}[h]
\centering
\caption{Magnitude of the change in cloud fraction across the VGT2 to Proba-V transition (cut at 2014-01-01) at the two sites of Fig.~3. Values are means of the daily cloud-fraction series, and the change is after minus before. Intervals are block-bootstrap 95\% intervals of the difference in means, with 30-day blocks and 5,000 replicates.}
\label{tab:discontinuity}
\small
\setlength{\tabcolsep}{5pt}
\begin{tabular}{llrrrrr}
\toprule
Site & Series & Days before & Days after & Mean before & Mean after & Change (95\% CI) \\
\midrule
SW Australia coast & Legacy (BS1/BS2) & 5\,204 & 2\,009 & 0.443 & 0.483 & $-0.040$ ($-0.068$ to $-0.016$) \\
SW Australia coast & Multi-FT PW-DL3   & 5\,204 & 2\,009 & 0.512 & 0.519 & $-0.007$ ($-0.036$ to $+0.018$) \\
BC glacier/snow     & Legacy (BS1/BS2) & 5\,421 & 2\,052 & 0.804 & 0.723 & $+0.080$ ($+0.051$ to $+0.112$) \\
BC glacier/snow     & Multi-FT PW-DL3   & 5\,414 & 2\,069 & 0.776 & 0.717 & $+0.060$ ($+0.029$ to $+0.094$) \\
\bottomrule
\end{tabular}

\vspace{6pt}
\parbox{0.95\textwidth}{\footnotesize The legacy record switches from BS1 to BS2 at the transition, so its change combines the sensor change with the change of algorithm. At the coastal site the change is about five times larger for the legacy record than for the proposed model. At the glacier site both records change, and the difference between them is smaller than at the coast.}
\end{table}
